\chapter{Fragmentos de los scripts ETL}
\label{ap:etl}

Este apéndice recoge tres fragmentos de código Python representativos del proceso ETL descrito en la Sección~\ref{sec:etl} y en la Sección~\ref{sec:metodologia_desarrollo}, elegidos por ilustrar patrones que se repiten a lo largo de todos los scripts de \texttt{scripts/} o por resolver una decisión de diseño no trivial. El código completo de cada script está disponible en el repositorio del proyecto.

% -----------------------------------------------------------------------
% A.1 PATRÓN DE IDEMPOTENCIA
% -----------------------------------------------------------------------
\section{Patrón de idempotencia}
\label{ap:sec:idempotencia}

Todos los scripts ETL comprueban, antes de hacer nada, si los datos que cargan ya están presentes en la base de datos, y solo actúan si no lo están o si la variable \texttt{OVERWRITE} se ha puesto explícitamente a \texttt{True} (Sección~\ref{sec:metodologia_desarrollo}). El Listado~\ref{lst:idempotencia}, tomado de \texttt{scripts/03\_load\_farolas.py}, muestra este patrón.

\begin{lstlisting}[style=Python, caption={Patrón de idempotencia (\texttt{scripts/03\_load\_farolas.py}).}, label={lst:idempotencia}]
with engine.connect() as conn:
    farolas_cargadas = conn.execute(
        text("SELECT count(*) FROM farolas")
    ).scalar() > 0

if farolas_cargadas and not OVERWRITE:
    print("Farolas ya cargadas en la BD, no se hace nada (OVERWRITE=False).")
else:
    if farolas_cargadas:
        with engine.begin() as conn:
            conn.execute(text("TRUNCATE TABLE farolas RESTART IDENTITY"))
    # ... extracción, transformación y carga de los datos ...
\end{lstlisting}

% -----------------------------------------------------------------------
% A.2 EXTRACCIÓN, TRANSFORMACIÓN Y CARGA CON GEOPANDAS
% -----------------------------------------------------------------------
\section{Extracción, transformación y carga con GeoPandas}
\label{ap:sec:etl_geopandas}

El Listado~\ref{lst:geopandas}, tomado de \texttt{scripts/06\_load\_distritos.py}, resume en pocas líneas las tres fases del proceso ETL descritas en la Sección~\ref{sec:etl}: lectura del fichero original en su formato nativo (\textit{extracción}), reproyección al sistema de referencia común EPSG:4326 y normalización de nombres de columna (\textit{transformación}), y volcado a PostGIS mediante \texttt{GeoPandas.to\_postgis()} (\textit{carga}).

\begin{lstlisting}[style=Python, caption={Extracción, transformación y carga de un Shapefile (\texttt{scripts/06\_load\_distritos.py}).}, label={lst:geopandas}]
shp_path = BASE / "data/raw/distritos/DISTRITOS.shp"
distritos = gpd.read_file(shp_path).to_crs("EPSG:4326")
distritos = distritos.rename(
    columns={"COD_DIS": "cod_distrito", "NOMBRE": "nombre",
             "Area": "area_m2", "geometry": "geom"}
)[["cod_distrito", "nombre", "area_m2", "geom"]]
distritos["cod_distrito"] = distritos["cod_distrito"].astype(int)
distritos = distritos.set_geometry("geom")

distritos.to_postgis("distritos", engine, if_exists="append", index=False)
\end{lstlisting}

% -----------------------------------------------------------------------
% A.3 MEDIA PONDERADA POR LONGITUD DE LOS INDICADORES DE UNA RUTA
% -----------------------------------------------------------------------
\section{Media ponderada por longitud de los indicadores de una ruta}
\label{ap:sec:media_ponderada}

Las estadísticas agregadas de cada ruta (peligrosidad media, iluminación media, vulnerabilidad media...) presentadas en el visor web (Sección~\ref{sec:frontend_visor}) y en la evaluación comparativa del Capítulo~\ref{cap:pruebas} no son una media simple de los tramos recorridos, sino una \textbf{media ponderada por la longitud de cada tramo}, para que un tramo de 200\,m no pese lo mismo que uno de 5\,m en el resultado agregado. El Listado~\ref{lst:media-ponderada}, tomado de \texttt{scripts/rutas.py}, implementa este cálculo.

\begin{lstlisting}[style=Python, caption={Cálculo de la media ponderada por longitud (\texttt{scripts/rutas.py}).}, label={lst:media-ponderada}]
distancia_m = sum(fila.length for fila in filas)

def media_ponderada(campo):
    return sum(
        fila.length * getattr(fila, campo) for fila in filas
    ) / distancia_m

peligrosidad_media = media_ponderada("indice_peligrosidad")
\end{lstlisting}
